\section{Boundaries and the next gate}
\label{sec:tpc47-boundaries}

The positive result is a transfer theorem for one canonical family of
orbit bands.  This section records every identification that would be
needed before it could enter a fixed-shift prime-pair argument.

\subsection{Five distinctions that remain active}

\paragraph{Residual synthesis versus the native affine face.}
At \(h_0\), the vector analyzed here is the actual-coefficient
negative large-prime residual synthesis \(-\sum_pUx_p\).  TPC--45 obtains this
object after averaging the lower prime signs.  Its primewise norms agree
with the corresponding native vectors, but its cross inner products do
not generally agree with those of the physical affine face
\citep{WangTPC45}.  The grouping statement in
\cref{prop:tpc47-grouped-component} recovers only the large-prime
component and pays the exact residual-fiber norm.

\paragraph{Frozen shifts versus reconstructed physical shifts.}
Only \(h=h_0\) has inherited physical provenance.  The shift transform
uses every integer \(h\) after freezing the \(h_0\) coefficient field.
For a newly reconstructed physical shift, the row set, residue factor,
mask, and source residues may change.  No growing-window uniformity
theorem is asserted.

\paragraph{Amplitude spectrum versus atomic normalization.}
Theorems \ref{thm:tpc47-exact-transfer} and
\ref{thm:tpc47-actual-transfer} concern the Fourier transform of a
Hilbert amplitude.  Fourier inversion at one \(h_0\) integrates all
frequencies, including the uncontrolled complement.  Moreover, the
TPC denominator is the atomic diagonal at \(h_0\), not the Cartesian
energy over all determinant shifts.  TPC--46 proves that these
normalizations may differ by a polynomial factor \citep{WangTPC46}.

\paragraph{Consecutive completion versus sparse aliases.}
The actual TPC completion samples determinant intercepts separated by
the triple-prime modulus.  Sparse sampling folds the completed shift
spectrum.  The present orbit projection neither unfolds those aliases
nor proves that the equality-alias residual \(d(m)r\) remains literal
at nonzero aliases; there the squarefree kernel can replace the product.

\paragraph{One band family versus the full orbit space.}
The complement \(I-\Pi_{\Omega}^{(H_0)}\) is uncontrolled.  By
\cref{thm:tpc47-no-mass-lower-bound}, qualitative injectivity of the
low sector supplies no quantitative reconstruction.  Summing separate
pointwise bounds over the \(X^{33/100+o(1)}\) comparable bands in
\eqref{eq:tpc47-cover-count} would introduce a prohibitive support
factor unless orthogonality is retained through the later operations.

\subsection{What has genuinely changed}

Before this paper, the orbit-output coupling was an arbitrary
\(j\)-dependent Hilbert vector, so the TPC--46 spectral theorem could
not be applied to any declared portion of the actual residual field.
Now there is a representation-independent norm-one operator, applied
to every output row and residual label, for which
\begin{enumerate}[label=\textup{(\roman*)},leftmargin=2.4em]
 \item the full actual coefficient reassembles with its source mask and
       M\"obius roles intact;
 \item a growing \(X^{1/400+o(1)}\) time--bandwidth sector has
       super-polynomial leakage on the interval
       \eqref{eq:tpc47-final-endpoint-gap}; and
 \item the exact operator-theoretic obstruction to uniform recovery is
       proved,
       rather than hidden in a phrase such as ``smooth output.''
\end{enumerate}
This crosses Door A of TPC--46 on one bounded projected residual sector.
It does not cross the other three doors.

\subsection{A reversible continuation}

The next paper may pursue any of the following logically separate
tasks.
\begin{enumerate}[label=\textbf{Door \Alph*.},leftmargin=3.7em]
 \item Partition the entire orbit circle into translated comb bands and
       prove a two-parameter resonance-tube square-function estimate
       that preserves orthogonality while the bands are reassembled.
       This avoids asking one low band to contain a fixed fraction of
       the vector.
 \item Use the literal TPC source/mask phases to prove a
       coefficient-specific spectral concentration or entropy bound.
       It must beat the arbitrary-vector counterexample in
       \cref{thm:tpc47-no-mass-lower-bound} by using arithmetic data.
 \item Combine orbit-band projection with the exact residual grouping
       map and prove that one additional cofactor range has acceptable
       grouping multiplicity at atomic normalization.
 \item Construct a family of actual masks or residual packets that
       saturates the high-frequency or grouping frontier.  This would
       give a sharp negative boundary without asserting failure of
       every possible TPC continuation.
\end{enumerate}

Door A is the most direct continuation: the projectors are orthogonal
before reassembly, and the bandwise resonance condition is explicit.
The missing theorem is a square-function or large-sieve estimate that
survives determinant dilation, residual grouping, and fixed-shift
localization.  No such estimate is assumed here.

\subsection{Prohibited conclusions}

Nothing in this paper proves or implies
\begin{enumerate}[label=\textup{(\roman*)},leftmargin=2.4em]
 \item deterministic cancellation of M\"obius coefficients at the
       prescribed physical shift;
 \item an atomic-normalized Bessel bound for the complete equality
       column;
 \item a uniform family of actual TPC packets over growing shifts;
 \item a parity-barrier breach, the Hardy--Littlewood prime-pair
       asymptotic, or the twin-prime conjecture.
\end{enumerate}
The advance is sectorial: the canonical low core has an exact gap, while
the projected frozen family built from the literal \(h_0\)-coefficient
field has arbitrary-power leakage there.  Its possible quantitative
thinness is also proved.
